\section{RAG y Curriculum: integrar la ubicación del conocimiento y el orden de aprendizaje en el scaling}

El capítulo anterior modificó el origen y la estructura lingüística del corpus; este capítulo cambia dos dimensiones adicionales: si el conocimiento se almacena en los parámetros o se recupera durante la inferencia, y si el modelo crece junto con un curriculum y un schedule de profundidad. Ambas rutas se oponen a la visión unidimensional de scaling basada únicamente en "ampliar el preentrenamiento estático".

\subsection{La interfaz básica de RAG: agregar memoria externa más allá de la memoria paramétrica}

Retrieval-Augmented Generation (RAG, generación aumentada con recuperación) recupera primero evidencias del repositorio de documentos según la consulta y luego entrega tanto la consulta como las evidencias al modelo de lenguaje. Desplaza parte de la actualización del conocimiento desde el entrenamiento paramétrico hacia la actualización de índices, pero incrementa la latencia de recuperación, los tokens de contexto, los errores de clasificación y las contradicciones en las evidencias.

\lecturefigure{slide-068.jpg}{Modelo base: una vez que los datos de preentrenamiento se comprimen en parámetros, generan la salida directamente}{CS25 V6 oficial deck, p.~68}

\lecturefigure{slide-070.jpg}{RAG completo: la consulta activa la recuperación y las evidencias entran al LLM como contexto}{CS25 V6 oficial deck, p.~70}

Si el presupuesto total de cómputo es $C$, se puede expresar aproximadamente como
\[
C=C_{\mathrm{pretrain}}+C_{\mathrm{index}}+C_{\mathrm{retrieve}}+C_{\mathrm{generate}}.
\]
$C_{\mathrm{pretrain}}$ forma el conocimiento paramétrico, $C_{\mathrm{index}}$ construye la base de datos externa, $C_{\mathrm{retrieve}}$ realiza búsquedas en cada solicitud y $C_{\mathrm{generate}}$ procesa el contexto aumentado. RAG no es conocimiento gratuito; redistribuye los costos desde un entrenamiento único hacia el almacenamiento y cada paso de inferencia.

\begin{knowledgebox}{Uso inicial: memoria paramétrica y no paramétrica}
La \textbf{memoria paramétrica} se refiere a las leyes estadísticas escritas en los pesos mediante gradientes; su actualización es lenta, la consulta es barata pero el origen no se puede localizar directamente. La \textbf{memoria no paramétrica} se refiere a repositorios de documentos, índices vectoriales o bases de datos; permite actualizaciones rápidas y rastreo del origen, pero depende de la calidad de la recuperación, los índices y el contexto adicional. RAG es una combinación de ambas, no "aprendizaje sin parámetros".
\end{knowledgebox}

\subsection{Scaling de RAG: el óptimo depende de la tarea y del presupuesto}

La investigación integra modelos base como OLMo-2 con módulos de recuperación en una asignación unificada de cómputo. Si hay demasiado poco preentrenamiento, el modelo podría no entender la consulta, las evidencias o la tarea de generación; si hay demasiado preentrenamiento, los beneficios marginales de continuar presionando hechos recuperables hacia los parámetros disminuyen. La proporción óptima de recuperación varía según la tarea: las problemas intensivos en conocimiento dependen más de las evidencias externas, mientras que el modelado lingüístico puro o la inferencia compleja siguen requiriendo un modelo base fuerte.

\lecturefigure{slide-071.jpg}{Estudiando el Scaling de RAG: escala paramétrica, tokens de preentrenamiento y recuperación cambian conjuntamente}{CS25 V6 oficial deck, p.~71}

\lecturefigure{slide-072.jpg}{¿Cuánto RAG es óptimo?: la asignación óptima de recuperación difiere según la tarea}{CS25 V6 oficial deck, p.~72}

\teachervoice{00:29:01--00:31:50, el docente enfatiza que existe un presupuesto mínimo útil de preentrenamiento; RAG no puede hacer que un modelo casi sin entrenamiento se vuelva fuerte automáticamente mediante la recuperación. El óptimo depende de la tarea, el recuperador y el cómputo total, no de una proporción fija.}

\begin{lstlisting}[language=Python,caption={Pseudocódigo para explorar la asignación entre preentrenamiento y recuperación bajo presupuesto fijo}]
for pretrain_budget in candidate_pretrain_budgets:
    model = train_base_model(tokens=budget_to_tokens(pretrain_budget))
    for retrieval_k in candidate_k:
        score = evaluate(
            model=model,
            retriever=index,
            k=retrieval_k,
            task=target_task,
        )
        record(pretrain_budget, retrieval_k, score)
choose_pareto_frontier()
\end{lstlisting}

\begin{warningbox}{Los tres costos de RAG más fáciles de ignorar}
En primer lugar, recuperar documentos aumentará los tokens de entrada y el caché KV; en segundo lugar, las evidencias erróneas empujarán al modelo hacia respuestas incorrectas pero "fundadas"; en tercer lugar, la actualización del índice y el control de acceso se convierten en nuevos estados del sistema. Comparar únicamente el número de parámetros sin contar los costos de recuperación/servicio llevará a malinterpretar el óptimo.
\end{warningbox}

En el lado de servicio también hay que distinguir entre prefill y decode. Las $k$ documentos recuperados se expanden una sola vez la longitud de entrada durante la fase de prefill, aumentando los costos de atención y establecimiento del caché KV; después, cada token de decode sigue leyendo estos cachés. Si las solicitudes de usuario son de alta concurrencia, un contexto de evidencia más grande reducirá el batch soportable. Por otro lado, un modelo base débil podría no utilizar los documentos correctamente, produciendo copiado de citas, contradicciones en las evidencias o ignorando la recuperación. Por ello, un experimento completo de scaling de RAG debe reportar al menos recall de recuperación, calidad de respuesta, longitud de contexto, latencia, throughput y costo de actualización del índice fresco, no solo la precisión final.

\subsection{Curriculum-Guided Layer Scaling: modelo y datos crecen juntos}

El aprendizaje por curriculum ordena las muestras según cierta dificultad o estructura. CGLS lleva esto más allá permitiendo que el conteo de capas crezca con el progreso del entrenamiento: el modelo primero procesa datos más simples con una red más superficial, e inserta gradualmente capas y aumenta la complejidad de los datos. Sea $s\in[0,1]$ el progreso de entrenamiento; el schedule de capas y dificultad se puede expresar como
\[
L(s)=L_0+\lfloor s(L_{\max}-L_0)\rfloor,
\qquad d(s)=d_0+s(d_{\max}-d_0).
\]
$L(s)$ es la profundidad actual y $d(s)$ es la métrica de dificultad de los datos. Los métodos prácticos necesitan definir cómo replicar/inicializar nuevas capas, cómo mantener el estado del optimizador y si la dificultad es realmente medible.

\lecturefigure{slide-075.jpg}{Motivación de CGLS: un modelo estático podría no coincidir con la evolución de la dificultad de los datos durante el entrenamiento}{CS25 V6 oficial deck, p.~75}

\lecturefigure{slide-076.jpg}{Motivación: cursos sintéticos, desarrollo infantil y DataComp-LM proporcionan pistas de orden}{CS25 V6 oficial deck, p.~76}

\lecturefigure{slide-077.jpg}{Trabajo previo: métodos de apilado de capas y crecimiento impulsado por similitud}{CS25 V6 oficial deck, p.~77}

\lecturefigure{slide-078.jpg}{Método CGLS: crecimiento de 8 a 16 capas según la fase de entrenamiento con cambios en el curriculum}{CS25 V6 oficial deck, p.~78}

\begin{knowledgebox}{Lectura de la figura: no interpretar "curso" solo como ordenamiento de muestras}
La diapositiva 78 cambia simultáneamente la profundidad del modelo, las capas transferidas, las capas inicializadas aleatoriamente y la complejidad de los datos. Antes de leer los resultados, pregunte qué parámetros se heredan, replican o crean en cada fase; de lo contrario, los cambios de rendimiento podrían provenir del crecimiento paramétrico, pasos de entrenamiento o inicialización, no solo del curriculum en sí.
\end{knowledgebox}

\begin{lstlisting}[language=Python,caption={Pseudocódigo para el schedule conjunto de profundidad del modelo y dificultad de datos}]
for stage in schedule:
    model = grow_layers(
        model,
        target_depth=stage.depth,
        transfer=stage.transfer_policy,
    )
    data = sample_by_difficulty(
        corpus,
        min_level=stage.min_difficulty,
        max_level=stage.max_difficulty,
    )
    train(model, data, steps=stage.steps)
\end{lstlisting}

\subsection{Resultados y límites: los pequeños beneficios requieren atribución correcta}

Las diapositivas 79--80 resumen múltiples benchmarks. Al comparar, primero identifique una línea base fija de cómputo/parámetros y luego determine si CGLS lidera consistentemente en el promedio, perplexity o tareas específicas. Incluso si existen beneficios promedio, verifique si sacrifica tareas individuales y si la complejidad de implementación adicional del proceso de crecimiento vale la pena.

\lecturefigure{slide-079.jpg}{Tabla de resultados de CGLS I: cambios promedio y por tarea en múltiples benchmarks}{CS25 V6 oficial deck, p.~79}

\lecturefigure{slide-080.jpg}{Tabla de resultados de CGLS II: comparación con diferentes líneas base y configuraciones}{CS25 V6 oficial deck, p.~80}

\lecturefigure{slide-081.jpg}{Discusión: beneficios, línea base normalizada y direcciones futuras del curriculum}{CS25 V6 oficial deck, p.~81}

\teachervoice{00:31:50--00:36:33, el docente describe CGLS como un método exploratorio: los resultados mejoran, pero aún se requiere investigación en la definición del curriculum, hiperparámetros, otras familias de modelos y escalas mayores. Experiencia práctica: no eleve una tabla de promedios a "leyes de crecimiento humanas".}

El crecimiento de capas también introduce problemas de estado del optimizador. Adam/AdamW mantienen un primer momento $m$ y un segundo momento $v$ para cada parámetro; cuando se agregan nuevas capas, estos estados del optimizador no tienen historia. Copiar capas existentes puede heredar representaciones, pero podría causar homogeneización entre capas; la inicialización aleatoria mantiene la diversidad, pero crea saltos distribucionales durante el entrenamiento. Las comparaciones justas deben explicar cómo se inicializan los pesos de las nuevas capas, $m/v$, el warmup de la tasa de aprendizaje y la escala residual, y emparejar la línea base según FLOPs totales o cantidad total de actualizaciones. De lo contrario, los beneficios de CGLS podrían provenir de etapas adicionales de entrenamiento en lugar del crecimiento conjunto del curriculum y la profundidad.

\subsection{Perspectiva unificada de datos: selección, mezcla, recuperación y curriculum}

Los cuatro proyectos se pueden resumir en una tabla de decisión: Baby Scale estudia la selección de fuentes, Bilingual BabyLM estudia la estructura de mezcla, RAG estudia la externalización del conocimiento y CGLS estudia el orden y el co-scaling. Juntos demuestran que la estrategia de datos no es solo filtrar antes del entrenamiento, sino una variable sistémica que atraviesa el aprendizaje paramétrico, la inferencia y el crecimiento del modelo.

\lecturefigure{slide-084.jpg}{Insight unificado I: origen de los datos, mezcla lingüística y dependencia de tareas}{CS25 V6 oficial deck, p.~84}

\lecturefigure{slide-085.jpg}{Insight unificado II: ubicaciones de intervención diferentes para curriculum, RAG y crecimiento del modelo}{CS25 V6 oficial deck, p.~85}

\lecturefigure{slide-086.jpg}{Conclusión general: los futuros modelos de lenguaje necesitan ser más inteligentes, no solo tener más datos}{CS25 V6 oficial deck, p.~86}

\begin{importantbox}{Lista de selección de estrategias de datos}
Si el problema es \textbf{ruido en el corpus de entrenamiento}, haga primero selección/filtrado; si el problema es \textbf{desequilibrio lingüístico o de dominio}, ajuste la mezcla; si el conocimiento se actualiza frecuentemente y requiere fuentes, considere RAG; si las etapas tempranas y tardías del entrenamiento requieren diferentes muestras, diseñe un curriculum; si la capacidad del modelo cambia según la fase, debe incluir el crecimiento de arquitectura en el experimento. Los diferentes controles no pueden reemplazarse con una sola métrica de "calidad de datos".
\end{importantbox}

\subsection{Resumen de la sección}

RAG extiende la ubicación del conocimiento desde los parámetros hacia la memoria externa, y CGLS extiende el orden de aprendizaje hacia un schedule conjunto de profundidad del modelo y dificultad de los datos. Ambos revelan que el núcleo del scaling es la asignación de recursos, no aumentar incondicionalmente tokens o parámetros. El siguiente capítulo gira hacia post-entrenamiento e inferencia de cómputo en tiempo real que pueden cambiar el comportamiento después de completar el entrenamiento del modelo.
